\section{Introduction}\label{sec:current-overview}\label{sec:introduction}

A predictive state contains the information in an experimental past that
can affect future laws. Its linear dimension measures one kind of
compression. A controller that stores probabilities on finitely many
labels must also respect positivity, the preparation of its initial
state, and the effect of subsequent observations. These constraints can
make its minimum memory much larger than the predictive dimension.
They also distinguish exact realization, uniform approximation of every
command word, and approximation of experiments stopped at bounded
expected cost. We study these distinctions within the original theory
of causal resource comparison and optimal statistical discrimination.

A terminal response assigns a probability row $b(w)$ to each finite
command word. An instrument response assigns joint probabilities to
command--report histories and includes the state on which a terminal
measurement acts. A classical realization uses one fixed collection
of nonnegative transition and readout tables. Its initial label is pure
and counted, except in explicitly named mixed-start comparisons.
The physical quantum system, an external policy's storage, program bits,
arithmetic precision and sampling work are separate resources. A private
simulator label is never made visible to the external policy.

The central question is how an algebraic or geometric description of
future responses controls their positive implementation. The present
results give complementary answers. A finite-group specification admits
a cost-preserving permutation reduction even when the required laws
have uncertainty or error. A centered pair of nonnegative coefficient
matrices has an explicit necessary and sufficient synthesis boundary,
which gives the complete mean--contrast resource law of an unbounded
nonabelian family. A return polytope and rare continuation branches give
finite approximations of entire stopped experiments, with memory governed
by the growth of the continuation semigroup.

\subsection*{Convex specifications and infinite-word approximation}

Fix a homomorphism $\gamma:A^*\to G$ onto a finite group. For each
$g\in G$ prescribe a closed convex set $K_g\subseteq\Delta(C)$ of
allowed terminal laws. The requirement is imposed on \emph{every} word
with image $g$. Theorem~\ref{thm:r76-convex-normal-form} proves that an
arbitrary stochastic controller on $M$ labels satisfies these constraints
if and only if there exist a finite $G$-set $X$, a preparation $\lambda$
and a stochastic decoder $D$ with
\[
 \lambda P_gD\in K_g\quad(g\in G),\qquad
 |X|+\mathbf1_{\{\lambda\text{ is mixed}\}}\le M.
\]
For a free mixed start the last condition is simply $|X|\le M$.
The original stochastic matrices are not assumed to satisfy group
relations, and the resulting action need not be faithful. The proof
averages identity words, projects onto closed communicating classes,
and shows that the sandwiched action permutes their stationary rows.
Closed convex constraints survive the limiting average. This extends
the retained exact normal form by replacing an exact observable
intertwining with a constraint-preserving projection.

Let $d_X(b)$ be the best uniform total-variation approximation to a
target $b:G\to\Delta(C)$ on the action $X$, and let
$d_X^{\mathrm{pt}}(b)$ restrict the preparation to a point.
Theorem~\ref{thm:r76-minimax-normal-form} gives the optimal error over
all words and all $M$-label pure-start controllers:
\begin{equation}\label{eq:r76-intro-minimax}
 \mathcal E_M^{\mathrm p}(b)
 =\min\left\{\min_{|X|\le M}d_X^{\mathrm{pt}}(b),
              \min_{|X|\le M-1}d_X(b)\right\}.
\end{equation}
Only finitely many actions occur. The minima are attained, depend
Lipschitz-continuously on the target, and are computable by real
quantifier elimination for algebraic data. This is synthesis for a
supplied finite-group target, not an algorithm for unrestricted
isolation or cutpoint questions about a supplied probabilistic automaton.

Theorem~\ref{thm:r76-robust-multiplicity} supplies lower certificates.
Fourier singular values larger than the allowed error force the
corresponding irreducible multiplicities in the core. Below half the
separation between distinct target residuals, any memory saving over
their number must also pay for a mixed initializer. For a prime cyclic
group, fewer than $p$ labels do no better than the best constant law,
whereas $p$ labels are exact (Corollary~\ref{cor:r76-prime-error}).
Approximation thus need not remove the algebraic memory obstruction
when the error must be uniform over words of unbounded length.

\subsection*{A sharp boundary with two positivity constraints}

The group average and the nonconstant response impose distinct
constraints on a probability factorization. For $n\ge5$ put $P=I-J/n$
and consider a probability column $u$ and a real matrix $A$ satisfying
\[
 0\le A\le u\mathbf1^{\mathsf T},\qquad
 \sum_{ij}A_{ij}=np,\qquad PAP=aP,\qquad
 \operatorname{rank}[A\mid u]<n.
\]
The matrices $A$ and $u\mathbf1^{\mathsf T}-A$ share a preparation
constraint. Theorem~\ref{thm:r76-bounded-coefficients} determines exactly
which $(p,a)$ permit such a pair, constructs a common nonnegative factor
of width $n-1$ at every permitted pair of parameters, and classifies the
coefficient pair on the curved boundary. This is an existence theorem;
it does not claim that every given feasible pair has that factor width.

Put
\[
 p_0=\frac1{n-\sqrt n},\qquad p_1=\frac2n,\qquad
 F_n(p)=\begin{cases}
 p(\sqrt n-1)/(\sqrt n+1),&0\le p\le p_0,\\
 ((n-1)p-1)/((n-1)(np-1)),&p_0\le p\le p_1,\\
 (1-p)/(n-1),&p_1\le p\le1.
 \end{cases}
\]
For $a>0$, the answer is $a\le F_n(p)$. The proof combines the retained
singular-matrix extremum with a simultaneous-positivity inequality.
It permits signed null vectors and nonsymmetric coefficient matrices.
The equality analysis forces the extremal coordinate pattern and then
the allowable common factors. On the middle interval both outcome
constraints are active; scaling the first-outcome bound alone would
give an incorrect threshold.

Let $\chi_n$ be the natural permutation character of $S_n$, and set
\[
 \beta_{p,a}(g)=p+a(\chi_n(g)-1),\qquad
 b_{p,a}=(\beta_{p,a},1-\beta_{p,a}).
\]
The exact probability range is
$0\le p-a\le1$ and $0\le p+(n-1)a\le1$.
For $a>0$, Theorem~\ref{thm:r76-mean-contrast} proves
\[
 \kappa_{S_n}(b_{p,a})=
 \begin{cases}n(n-1),&a\le F_n(p),\\n^2,&a>F_n(p),\end{cases}
 \qquad m(b_{p,a})=1+\kappa_{S_n}(b_{p,a}).
\]
Here $\kappa$ is the least positive permutation core and $m$ counts a
pure start. Negative contrast exchanges the outcomes and replaces
$p$ by $1-p$; zero contrast has one-label memory. At every nonzero
contrast the observable group is $S_n$, the residual number is $n!$,
and the predictive dimension is $1+(n-1)^2$. The mean $p$ is an
average response over the group, not a prior on an unknown parameter.
All minimum cores on the low and middle transition boundary are
classified. The middle boundary forces equal orbit weights, while
the low-mean boundary permits specified nonuniform weights.
For $p\ge2/n$, the lower resource suffices throughout the positive-contrast
probability range. Two commands and binary readout realize the family
as a prepared finite-dimensional quantum experiment.

The retained arbitrary-group results explain its representation-theoretic
part: small-positive-contrast memory equals a permutation-module
multiplicity cover (Theorem~\ref{thm:r75-orbit-multiplicity}), and a
single transitive response set is convex exactly at normal stabilizers
(Theorem~\ref{thm:r75-normality-convexity}). These proofs, the exact
R75 return law, the quotient linear programs and their payoff duals
remain current inputs in the supplement.

\subsection*{Complete stopped laws and finite return geometry}

Each observed instrument report changes the conditional state used by
later calls. Theorem~\ref{thm:r76-return-approximation} assumes a finite
physical polytope containing the preparation. Every returning branch
maps the invariant state set into that polytope; continuing branches
have total conditional probability at most $r<1$. Retain all distinct
normalized continuation images of the seed vertices of length less
than $L$, and denote their number by $G(L)$.

One $G(L)$-label mixed-start model, or a counted pure-start model with
at most $G(L)+1$ labels, is exact before the first truncated excursion.
Uniformly over visible adaptive command, stopping and terminal-POVM
policies of target expected calls at most $B$, its complete stopped-law
error is at most
\begin{equation}\label{eq:r76-intro-excursion}
 \min\left\{1,\frac{(1-r)B}{r^{-L}-1}\right\}.
\end{equation}
The bound also holds for the terminal classical--quantum record,
including nontermination mass. The exceptional-run estimate is sharp
as a stopping-probability inequality
(Lemma~\ref{lem:r76-excursion-budget}); this sharpness alone does not
give a lower bound for every approximating controller.

With $b$ commuting continuation maps and $k$ seeds,
$G(L)\le k\binom{L+b-1}{b}$. Other polynomial orbit-growth estimates
give the corresponding power of $\log(B/\delta)$ for accuracy $\delta$.
The hypotheses concern individual normalized branches, and are
verified directly in the application; average mixing is not
substituted for a return condition.

For the retained undamped three-report qubit instrument, a regular
triangle in the real Bloch plane contains every report-$2$ posterior;
report $1$ is a pure reset, and report $0$ is a rotation of probability
$1/8$. For any fixed known alphabet of $a$ commanded angles,
Corollary~\ref{cor:r76-qubit-polylog} gives explicit branch tables with
\[
 M\le1+3\binom{L+a-1}{a},\qquad
 L=\max\left\{1,\left\lceil
       \log_8\left(1+\frac{7B}{8\delta}\right)\right\rceil\right\}.
\]
One angle therefore needs at most $3L+1$ labels, uniformly in its
numerical value. This improves the damping upper for fixed finite
alphabets. The earlier $O(\sqrt{B/\delta})$ construction still applies
simultaneously to the full angle circle. These are different command
quantifiers, not competing assertions about one unrestricted model.

Within the class of fixed physical-state dictionaries with Born terminal
decoders, the one-angle upper has a matching lower. For a fixed badly
approximable rotation phase and fixed $B>0$,
Theorem~\ref{thm:r76-physical-logarithm} proves
\[
 \lim_{\delta\downarrow0}
 \frac{N_{\rm phys}^{\rm m}(\delta;B)}{\log_8(1/\delta)}
 =\lim_{\delta\downarrow0}
 \frac{N_{\rm phys}^{\rm p}(\delta;B)}{\log_8(1/\delta)}=1.
\]
Rare reset--rotation words prepare separated pure states. Joint stopped
record events force distinct physical dictionary points to approximate
those states; zero-call readouts exclude the pure initializer from that
packing. An orbit-prefix return polygon attains the leading constant.
The golden-ratio phase has an explicit arithmetic certificate and five
seed states. The original algebraic rotation has a rational-table upper
with seven seeds, without an asserted badly approximable certificate.
This lower is deliberately stated for physical dictionaries; arbitrary
classical terminal decoders need not admit that representation.

The general stationary-initialization identity remains in the main
article: except for the one-label case, a predictively stationary
response costs exactly one label beyond its least mixed-start
realization (Theorem~\ref{thm:r75-stationary-initializer}). The complete
noisy-instrument laws, exact sector and worst-phase counts, all-angle
dictionary and boundary rigidity remain current in the supplement.
Approximate truncation need not be stationary, so its extra initializer
is an upper construction rather than an automatic lower bound.

\subsection*{Priority, dependencies, and the original causal programme}

Residual equivalence and stochastic automata are classical
\cite{RabinScott59,Rabin63,Heller65,CarlylePaz71}. An exact deterministic
residual theorem is not itself a new positive-realization principle.
Rabin's isolated-cutpoint reduction and approximate matrix-product
questions have different quantifiers from
\eqref{eq:r76-intro-minimax}. Nonnegative stochastic groups
\cite{Schwarz64,Flor69,Brown64,Schwarz62}, finite-group representations
\cite{Serre77}, equivariant factorization and nonnegative-rank geometry
\cite{GPT13,KPT10,BocciCarliniRapallo11,GillisGlineur12} provide
established ingredients. Bounded simplex factorization and shared
complementary-output factors are also prior theory
\cite{VuThanhGillisLecron23}; generic nonnegative-rank thresholds and
nested-polytope transitions have earlier examples \cite{Gillis12}. Real-algebraic minimization of probabilistic
automata is also prior work \cite{KieferWachter14}; the finite-group
projection makes the entire wordwise approximation constraint finite.

Polyhedral simulation and adaptive composition of quantum instruments
have explicit antecedents \cite{MonrasWinter14,MonrasWinter16,Okay26}.
Phase obstructions and noise-dependent classical-memory lower bounds
are treated by Fanizza, Lumbreras and Winter \cite{Fanizza24}.
Waiting-time martingales for patterns are classical \cite{Li80}.
The added claims are the stated convex-specification reduction, the
simultaneous-positivity synthesis boundary, and the complete adaptive
stopped-law construction from return geometry. A targeted literature
search and finite diagnostics do not certify exhaustive priority.

The original actual-history comparison theorem and its common paid
construction remain in the main article. Its witness is evaluated on
the competitor's actual pre-call histories. The paired supplement keeps
the complete arguments on performed-law compactness, variational and
dual attainment, robust tied optimality, stopping rigidity, finite-width
transfer, rank-constrained discrimination, finite-order quantum
measurement, two-stage reset control, memory and communication,
spectral geometry and optimal strategy implementation. The independent
R58 companion retains calibration and nonintervening directional
sensing. Local geometric conclusions keep their local hypotheses;
independent A1, A2 and further Theta objectives remain separate
research problems. The theorem map distinguishes logical prerequisites
of the current article from historical preservation.
